\section{Economic position}
\label{sec:economics}

\subsection{No rent on the medium of exchange}

\begin{property}[Absence of monetary rent]
\label{prop:no-rent}
No party earns income from the operation of the medium. There is no fee token, no gas,
no seigniorage, no interest on capacity, no treasury and no protocol revenue. Validators
are community institutions bearing a cost, not intermediaries extracting one.
\end{property}

A structural commitment rather than a policy setting, and it constrains the rest of the
design: spam cannot be priced by fees, so it is priced by refundable bonds
(\S\ref{sec:recourse}); capacity cannot be sold, so it must be earned; loss cannot be
absorbed by a treasury that does not exist, so mutualization rests on a standing, signed
liability somebody accepted in advance --- a supply declaration, the only thing in the
model that can carry a loss.

\subsection{Bounded standing, and nothing to hoard}

\begin{property}[No hoardable store of value]
\label{prop:no-hoard}
The protocol has no balance. What a member accumulates is capacity, which is
non-transferable, non-purchasable, bounded by a cut, decaying unless renewed, and held
in other members' stakes rather than in the member's own account.
\end{property}

A member cannot accumulate the medium, corner it, lend it at interest or exit into it. A
large claim is a large expectation of future delivery, and it decays into nothing if the
community around it stops trading: what one accumulates is a relationship, and
relationships lapse. Capacity concentrates only to the extent that many members
individually choose to stake on the same account, and even then
Corollary~\ref{cor:uw-bound} bounds the result; no mechanism lets holding capacity
generate more.

\subsection{Loss, mutualized but bounded}

Where a loss falls is decided by the tier, not by an ordering between remedies. An
\textbf{insured} loss falls on the underwriters whose supply carried it, split by exactly
what each arc carried and bounded by Corollary~\ref{cor:uw-bound}; if one of them defaults
in turn, on the creditor, visibly. An \textbf{uninsured} loss falls on the creditor who
chose it, alone (Remark~\ref{rem:uninsured-alone}). There is nothing after that, because
there is nothing else to draw on: no lender of last resort. A community that could
socialize unbounded loss would have created an unfunded liability on its members, the
failure mode this design most wants to avoid, since it converts a credit system into a
claim on people who never agreed to underwrite it.

\subsection{Autonomy and regulatory scope}

Obligations here are commercial trade credit between counterparties who have transacted,
denominated in a community unit, non-transferable to third parties as an instrument, and
not redeemable for currency. The medium is not a deposit, not a security and not a payment
token; it is the bookkeeping of mutual trade credit, made explicit. That describes the
mechanism, is not legal advice, and varies by jurisdiction. What the design provides for is
auditability: a community can publish its state root, prove any member's position, and
export a signed history.

\subsection{What conventional money still does better}

Money is accepted by strangers; capacity here is worth nothing to a community one has not
traded with. Money stores value across decades; capacity decays within epochs. Money is
anonymous at the point of sale; obligations here are between identified counterparties
within a community. Money can be saved by an individual acting alone; provision here runs
through named counterparties, and at long horizons through institutions
(\S\ref{sec:provision}). This design does not replace money. The claim is narrower: for
the subset of economic life that consists of regular mutual trade among producers who know
one another, a clearing medium fits better than a scarce asset, and that subset is larger
than it is usually taken to be.
